\clearpage\fancyhead[L]{\small KOMET.}\hypertarget{native-347-lower}{}\label{native-347-lower}
\begin{center}\textbf{KOMET.}\par\vspace{4pt}\hrule\vspace{4pt}\fontsize{10}{12}\selectfont Göttingische gelehrte Anzeigen. Heft 201, S. 2001\,..\,2004, 19. Dezember 1811.\par\vspace{4pt}\hrule\vspace{4pt}\end{center}

\par\indent In Heft 151 dieser Blätten [S. 337 d. B.] haben wir die ersten hiesigen Beobachtungen des großen diesjährigen Kometen zusammen mit den ersten verbesserten parabolischen Elementen bekannt gegeben. Diese Elemente waren lediglich als vorläufig bezeichnet worden. Sie wichen bald um einige Minuten von den Beobachtungen ab, und Anfang Oktober war der Unterschied bereits auf einen Viertelgrad angewachsen. Eine leichte Verbesserung der parabolischen Elemente genügte jedoch, um diesen Fehler zu beseitigen. Diese zweiten verbesserten Elemente wurden im Oktoberheft der Monatl. Correspondenz [S. 338 d. B.] bekannt gemacht. Anfang November war der Fehler dieser Elemente, deren Berechnung nur einige einzelne Beobachtungen zugrunde gelegen hatten, wieder auf einige Minuten angewachsen. Herr Professor \textsc{Gauß} hielt es daher für interessant, zu untersuchen, inwieweit sich dies schon als Beweis für die Elliptizität der Bahn ansehen ließe. Er ließ die dafür nötigen Rechnungen unter seiner Aufsicht von Herrn \textsc{Nicolai} ausführen. Von dessen hervorragender Geschicklichkeit und Sorgfalt beim astronomischen Rechnen haben wir schon früher in diesen Blättern Beispiele mitgeteilt. Man legte den größten Teil der Beobachtungen zugrunde, die Herr von \textsc{Zach} in der ersten Sichtbarkeitsperiode des Kometen gemacht hatte, sowie eine große Zahl neuerer Beobachtungen, die bis zum 6. November reichten. Aus dem Vergleich dieser Beobachtungen mit den letzten Elementen von Herrn Professor \textsc{Gauß} leitete man vier Normalorte ab. Es zeigte sich, dass sich die neueren Beobachtungen mit einer nur sehr kleinen Korrektur dieser letzten Elemente genau darstellen ließen, während bei den älteren Beobachtungen nur kleine Unterschiede blieben, nämlich $16''$ in der Länge und $28''$ in der Breite beim ersten Normalort vom 16. April sowie $51''$ in der Länge und $120''$ in der Breite beim zweiten vom 18. Mai. Zwar ist kaum anzunehmen, dass der letztere Normalort, der Mittelwert aus einer großen Zahl freilich nicht sehr genauer Beobachtungen, tatsächlich um den ganzen Betrag dieses großen Unterschieds fehlerhaft ist. Doch solange ungewiss ist, wie viel davon auf den Normalort selbst entfällt, ist der Unterschied noch zu klein, um darauf eine einigermaßen zuverlässige Bestimmung der Ellipse zu gründen. Herr Professor Gauss hielt es daher für besser, diese Arbeit zu verschieben, bis sich anhand späterer Beobachtungen etwas entscheiden lässt. Auf jeden Fall liegt die Umlaufzeit weit über 1000 Jahren. Die verbesserten parabolischen Elemente nach der Berechnung von Herrn \textsc{Nicolai} lauten wie folgt.

\Needspace{18\baselineskip}
\begin{center}\fontsize{10}{12}\selectfont\setlength{\tabcolsep}{4pt}\renewcommand{\arraystretch}{1.15}\begin{tabular}{l}
Durchgang durch das Perihel am 12. September 1811\\
\quad $6^{\mathrm u}\;30^{\mathrm m}\;55^{\mathrm s}$ Mittl. Z. in Göttingen\\
\end{tabular}\end{center}
\begin{center}\fontsize{10}{12}\selectfont\setlength{\tabcolsep}{3pt}\renewcommand{\arraystretch}{1.15}\begin{tabular}{llr}
Logarithmus des kleinsten Abstands &  & $0.0151048$\\
Länge des Perihels & .\;.\;.\;.\;. & $75^\circ\;1'\;44''3$\\
Länge des aufsteigenden Knotens & .\;.\;. & $140\;21\;57.5$\\
\end{tabular}\end{center}
\noindent Beide sind gegenüber den Fixsternen fest und werden vom Äquinoktium des 12. September aus gezählt.
\begin{center}\fontsize{10}{12}\selectfont\setlength{\tabcolsep}{3pt}\renewcommand{\arraystretch}{1.15}\begin{tabular}{llr}
Neigung der Bahn & .\;.\;.\;.\;.\;.\;. & $73^\circ\;4'\;30''9$\\
\end{tabular}\end{center}
\noindent Bewegung rückläufig.

\SourceLeaf{348}
\par\indent Um die Beobachtungen im Januar des nächsten Jahres zu erleichtern, hat Herr \textsc{Nicolai} noch die folgende Ephemeride berechnet. Wir geben sie hier vollständig wieder, auch wenn der Komet schon vorher nicht mehr sichtbar sein wird. \mbox{--- --- ---}


\par\indent Das Vorstehende war bereits für den Druck niedergeschrieben, als Herr Professor \textsc{Gauß} in einem Brief von Herrn von \textsc{Lindenau} erfuhr, dass Herr \textsc{Pons} am 16. November in Marseille einen weiteren neuen Kometen im Eridanus entdeckt hatte. Dem Brief lagen die folgenden Beobachtungen von Herrn von \textsc{Zach} in Marseille bei.

\begin{center}
\fontsize{10}{12}\selectfont\setlength{\tabcolsep}{3pt}\renewcommand{\arraystretch}{1.15}\begin{tabular}{l|c|c|c}
1811 & \shortstack{Mittlere Zeit\\in Marseille} & \shortstack{Rektaszension} & \shortstack{Südliche\\Deklination} \\ \hline
Nov. 17 & $10^{\mathrm u}\;23^{\mathrm m}$ & $67^\circ.25'$ & $25^\circ\,52'$ \\
18 & $11\;\;11\;\;17^{\mathrm s}3$ & $67\;\;14\;\;39''8$ & $25\;\;24\;\;\;8''6$ \\
19 & $\;9\;\;59\;\;32.8$ & $67\;\;\;4\;\;59.6$ & $24\;\;54\;\;\;8.5$ \\
20 & $10\;\;\;8\;\;37.7$ & $.66\;\;56\;\;\;8.2$ & $24\;\;18\;\;\;9.2$ \\
21 & $10\;\;14\;\;45.5$ & $66\;\;46\;\;53.0$ & $23\;\;41\;\;47.8$ \\
\end{tabular}
\end{center}


\Needspace{22\baselineskip}\par\indent Gleich am Abend des 9. Dezember, an dem diese Nachricht eingegangen war, begünstigte ein sehr klarer Himmel die Suche nach dem Kometen an der hiesigen Sternwarte. Herr Professor \textsc{Harding} sichtete ihn sofort mit einem Kometensucher in der Nähe eines Sterns siebter Größe. Herr Professor \textsc{Gauß} verglich den Kometen mehrmals mit diesem Stern am Kreismikrometer. Die scheinbare Position des Sterns wurde aus der \emph{Histoire céleste} zu $64^\circ\,11'\,23''1$ Rektaszension und $10^\circ\,22'\,2''5$ südlicher Deklination bestimmt. Daraus ergab sich die folgende Position des Kometen.

\Needspace{9\baselineskip}
\begin{quote}\raggedright 1811, 9. Dez. $10^{\mathrm u}\;6^{\mathrm m}\;52^{\mathrm s}$ Rektasz. $63^\circ\;49'\;41''4$, südliche Deklination $10^\circ\;21'\;55''5$\par\medskip Am 11. Dez. um $10^{\mathrm u}\;34^{\mathrm m}\;1^{\mathrm s}$ M. Z. ergab sich eine Rektasz. von $63^\circ\;33'\;20''5$ und eine südl. Dekl. von $8^\circ\;39'\;54''7$\end{quote}


\par\indent Dieser neue Komet ist bisher übrigens sehr klein und lichtschwach, etwa so hell wie ein Stern achter Größe. Er scheint sich jedoch der Erde zu nähern, da sich seine Bewegung beschleunigt. Weil er einen deutlicheren Mittelpunkt als der große Komet zeigt, lässt er sich am Kreismikrometer auch besser beobachten.


\par\indent Am 12. Dezember.


